\begin{recipe}{Lekkerbekjes, aardappelkroketjes en sla}
\kicker[Hoofdgerecht,Vis]{Hoofdgerecht · vis}
\index[register]{Lekkerbekjes, aardappelkroketjes en sla}
\index[register]{Kibbeling}
\index[register]{Aardappelkroketjes}
\index[register]{IJsbergsla}

\meta{20 minuten}

\lettrine{D}{e} kinderen noemen dit gerecht naar onze dochter, zo lekker vindt ze
het. Het is ook een van de weinige visgerechten die de kinderen lusten.
Kant-en-klare kibbeling en kroketjes uit de oven, met een simpele frisse sla
erbij. Niet veel werk, wel een vaste favoriet.

\blockrule{Ingrediënten}
\begin{ingredients}
  \ing{750 g}{kibbeling (Visgilde)}\index[register]{Kibbeling}
  \ing{600 g}{aardappelkroketjes uit de oven}\index[register]{Aardappelkroketjes}
  \ing{1 krop}{ijsbergsla}\index[register]{IJsbergsla}
  \ing{1}{tomaat}
  \ing{1}{komkommer}
  \ing{1 zakje}{slapitjes (Lidl)}
  \ing{3 el}{olijfolie}
  \ing{1 el}{balsamicoazijn}
  \ingb{remouladesaus}
  \ingb{augurken of zilveruitjes}
  \ingb{versgemalen peper \& zout}
  \ingb{verse dille}
\end{ingredients}

\blockrule{Bereiding}
\begin{steps}
  \item Verwarm de oven vast voor volgens de aanwijzingen op de verpakking van
        de kroketjes, dan is hij op temperatuur zodra je ze erin zet.
  \item Haal de kibbeling op; die blijft vanzelf nog wel even warm.
  \item Bereid de aardappelkroketjes in de oven volgens de aanwijzingen op de
        verpakking. Is de kibbeling toch wat afgekoeld, zet hem dan de laatste
        5 minuten in een ovenschaal mee met de kroketjes om weer op te warmen.
  \item Snijd ondertussen de ijsbergsla, de tomaat en de komkommer en doe ze in
        een ruime schaal. Voeg de slapitjes toe.
  \item Meng de olijfolie met de balsamicoazijn, peper, zout en dille tot een
        dressing en schep die door de sla.
  \item Serveer de kibbeling en kroketjes met de sla erbij.
\end{steps}
\end{recipe}
